\section{Global polynomial and two-branch rigidity}\label{sec:rigidity}
Write an ordinary periodic orbit as $(y_i,y_{i+1})$, with indices
cyclically extended. Its recurrence is
\begin{equation}\label{eq:recurrence}
 y_i^2+c=y_{i+1}+a y_{i-1}.
\end{equation}
The inverse map is $((x^2+c-y)/a,x)$, so this recurrence is valid
without excluding periods one and two or inserting a preperiodic tail.
Polynomial degree always means degree in $t$, with $\deg0=-\infty$
in degree inequalities.

\begin{lemma}[Finite-pole escape]\label{lem:integrality}
Every rational periodic coordinate lies in $k[t]$.
\end{lemma}
\begin{proof}
For an irreducible polynomial $\pi\in k[t]$, let $v_\pi(\pi)=1$.
If some coordinate has a pole at $\pi$, choose one of maximal pole
order $M>0$ within its finite orbit. At that index the left side of
\eqref{eq:recurrence} has valuation exactly $-2M$, because $c$ is
polynomial. The right side has valuation at least $-M$, because
$a$ is a nonzero constant. This is impossible. Thus no coordinate
has a finite polynomial pole. A coprime rational fraction with a
nonconstant denominator would have such a pole, proving the claim.
\end{proof}

\begin{lemma}[Common degree across all cycles]\label{lem:degree}
If a periodic point exists, $\deg c=2m$ for some $m>0$, and every
coordinate in every periodic orbit has degree $m$.
\end{lemma}
\begin{proof}
For one orbit let $m$ be the largest coordinate degree. If $m\le0$,
all coordinates are constant, contrary to \eqref{eq:recurrence}
and the nonconstancy of $c$. At a coordinate of degree $m>0$, the
right side of the recurrence has degree at most $m$. If $\deg c$
were smaller or larger than $2m$, the left side would have degree
$2m$ or $\deg c$, respectively. Therefore $\deg c=2m$ and the
leading terms cancel. A coordinate of degree below $m$ could not
cancel the leading term of $c$, giving the same contradiction.
The common value $m=\deg(c)/2$ depends only on the parameter,
so it is the same in every orbit. This argument does not take a
maximum over a periodic set not yet known to be finite.
\end{proof}

\begin{proposition}[Centered parameter and global branches]\label{prop:branches}
If a periodic point exists, \eqref{eq:centered} holds with the
stated uniqueness. With $P$ fixed, every periodic coordinate is
uniquely $\sigma P+b$ for $\sigma\in\{1,-1\}$ and $b\in k$.
\end{proposition}
\begin{proof}
Choose one coordinate $P_0$ in one orbit. If $y$ is a coordinate in
any periodic orbit, subtract the recurrence at $P_0$ from that at
$y$. Lemma~\ref{lem:degree} gives
\begin{equation}\label{eq:square-difference}
 \deg(y^2-P_0^2)\le m.
\end{equation}
Both coordinates have degree $m$. Their leading coefficient ratio
has square one, hence equals a unique $\sigma\in\{1,-1\}$.
Since $2\ne0$, the polynomial $y+\sigma P_0$ has degree $m$.
Factoring the left side of \eqref{eq:square-difference} shows that
$y-\sigma P_0$ is constant, also allowing it to be zero.
This holds across all cycles.

Apply the recurrence at $P_0$. Its two neighbors now show that
$c=-P_0^2+U P_0+V$ for some $U,V\in k$. Set
$P=P_0-U/2$ and $C=V+U^2/4$. Then $c=-P^2+C$, and the
constant change in $P_0$ leaves every coordinate of the asserted
form. Its sign and constant are unique because $P$ is nonconstant.

If also $c=-Q^2+D$ with nonconstant $Q$ and constant $D$, then
$(P-Q)(P+Q)=C-D$. If $C-D\ne0$, both factors are nonzero
constants, which would make $2P$ constant. Hence $C=D$, and
the integral domain $k[t]$ gives $Q=P$ or $Q=-P$.
This is a normalization of the parameter, not a change to the map.
\end{proof}
